\section{Lean y la demostración formal de teoremas}

\subsection{Introducción a Lean}

Lean es un lenguaje de programación, un demostrador de teoremas y un asistente de demostración interactivo:
\begin{itemize}
    \item Define conceptos básicos como los números naturales y la suma
    \item Declara y demuestra teoremas (como la conmutatividad $a+b=b+a$)
    \item Una demostración se representa como un \textbf{árbol de demostración}: el nodo raíz es el teorema original, y cada paso lo descompone en subobjetivos más simples
    \item Los archivos se organizan en proyectos, y estos pueden reutilizarse entre sí (de forma similar a las dependencias de bibliotecas en el desarrollo de software)
\end{itemize}

\subsection{LeanDojo: infraestructura abierta para la demostración de teoremas}

\begin{importantbox}{LeanDojo (NeurIPS 2023)}
Proporciona conjuntos de datos, herramientas y modelos de código abierto:
\begin{itemize}
    \item Extrae alrededor de 10,000 teoremas y demostraciones a partir de código Lean escrito por humanos
    \item Extrae el estado de la demostración en cada paso, los lemas utilizados y sus definiciones
    \item Entrena \textbf{ReProver} (demostrador aumentado por recuperación): primero recupera los lemas relevantes y luego genera el siguiente paso
\end{itemize}
\end{importantbox}

\subsection{Mejora mediante Expert Iteration}

Se usa el demostrador actual para intentar demostrar nuevos teoremas $\to$ se recopilan las demostraciones exitosas $\to$ se agregan al conjunto de entrenamiento $\to$ se entrena un demostrador más fuerte $\to$ se itera, hasta que el desempeño se satura.

\subsection{El desafío fundamental de la demostración de teoremas}

\begin{warningbox}{Espacio de acciones infinito}
El go tiene un tablero finito de $19\times19$, pero el espacio de acciones de una demostración matemática es \textbf{prácticamente infinito}: es difícil cubrirlo solo con datos humanos, y la exploración con RL también es extremadamente difícil.
\end{warningbox}

\subsection{Caso: demostración de desigualdades (LiPS)}

En este dominio específico de las desigualdades de las olimpiadas de matemáticas:
\begin{itemize}
    \item Las acciones se dividen en \textbf{escalado} (aplicar lemas conocidos, finito y enumerable) y \textbf{reescritura} (transformar la fórmula, un espacio infinito que se deja al LLM)
    \item Un algoritmo simbólico enumera todos los escalados + el LLM se encarga de la reescritura + filtrado heurístico
    \item O1 Preview demuestra 0/20, O3 y DeepSeek-R1 demuestran cada uno 3--4/20, los medallistas de oro humanos 15/20, \textbf{LiPS demuestra 16/20}
    \item Se descubrió un método de demostración desconocido para los humanos (usando solo AM-GM junto con una reescritura ingeniosa)
\end{itemize}
